\textbf{Phần 1. Trắc nghiệm lựa chọn}
\setcounterpageref{ex}{0}

\begin{ex} %[3X15N75]
Dùng các vạch kiểm đếm: mỗi nhóm gồm 4 gạch dọc và 1 gạch chéo biểu thị số lượng là
\choice
{\True $5$}
{$4$}
{$6$}
{$10$}
\loigiai{
Mỗi nhóm gồm 5 đơn vị.\\
\textbf{Đáp án của câu hỏi là A}
}
\end{ex}

\begin{ex} %[3X15N75]
Có 3 nhóm kiểm đếm (mỗi nhóm 5 gạch) và 2 gạch lẻ biểu thị số lượng là
\choice
{$15$}
{\True $17$}
{$16$}
{$18$}
\loigiai{
$5 \times 3 + 2 = 17$.\\
\textbf{Đáp án của câu hỏi là B}
}
\end{ex}

\begin{ex} %[3X15H75]
Thực hành tung đồng xu 20 lần: xuất hiện 12 lần mặt sấp. Số lần xuất hiện mặt ngửa là
\choice
{\True $8$ lần}
{$7$ lần}
{$9$ lần}
{$10$ lần}
\loigiai{
$20 - 12 = 8$ lần.\\
\textbf{Đáp án của câu hỏi là A}
}
\end{ex}

\begin{ex} %[3X15H75]
Thực hành gieo xúc xắc 30 lần: mặt 6 chấm xuất hiện 6 lần. Mặt 6 chấm chiếm bao nhiêu phần tổng số lần gieo?
\choice
{$\dfrac{1}{4}$}
{$\dfrac{1}{6}$}
{\True $\dfrac{1}{5}$ (vì $30 : 6 = 5$)}
{$\dfrac{1}{3}$}
\loigiai{
$30 : 6 = 5$ nên chiếm $\dfrac{1}{5}$.\\
\textbf{Đáp án của câu hỏi là C}
}
\end{ex}

\begin{ex} %[3X15H75]
Kiểm đếm phương tiện giao thông đi qua cổng trường trong 15 phút: Xe đạp: 25 xe, Xe máy: 40 xe, Ô tô: 10 xe. Phương tiện nào đi qua nhiều nhất?
\choice
{Xe đạp}
{Ô tô}
{Xe đạp và xe máy bằng nhau}
{\True Xe máy}
\loigiai{
Xe máy nhiều nhất (40 xe).\\
\textbf{Đáp án của câu hỏi là D}
}
\end{ex}

\begin{ex} %[3X15H75]
Tổng số phương tiện đi qua cổng trường trong thời gian trên là
\choice
{$70$ xe}
{\True $75$ xe}
{$80$ xe}
{$85$ xe}
\loigiai{
$25 + 40 + 10 = 75$ xe.\\
\textbf{Đáp án của câu hỏi là B}
}
\end{ex}

\begin{ex} %[3X15V75]
Khảo sát món ăn sáng yêu thích của 40 bạn học sinh: Bánh mì: 15 bạn, Phở: 12 bạn, Xôi: 8 bạn, còn lại thích Bún. Số bạn thích Bún là
\choice
{$4$ bạn}
{$6$ bạn}
{\True $5$ bạn}
{$7$ bạn}
\loigiai{
$40 - (15 + 12 + 8) = 40 - 35 = 5$ bạn.\\
\textbf{Đáp án của câu hỏi là C}
}
\end{ex}

\begin{ex} %[3X15V75]
Món ăn sáng được nhiều bạn yêu thích nhất là
\choice
{Phở}
{Xôi}
{Bún}
{\True Bánh mì}
\loigiai{
Bánh mì được thích nhất (15 bạn).\\
\textbf{Đáp án của câu hỏi là D}
}
\end{ex}

\begin{ex} %[3X15V75]
Lấy ngẫu nhiên một quả bóng trong hộp 15 lần (có hoàn lại): bóng đỏ xuất hiện 9 lần, bóng vàng xuất hiện 6 lần. Bóng đỏ xuất hiện nhiều hơn bóng vàng số lần là
\choice
{$2$ lần}
{\True $3$ lần}
{$4$ lần}
{$5$ lần}
\loigiai{
$9 - 6 = 3$ lần.\\
\textbf{Đáp án của câu hỏi là B}
}
\end{ex}

\begin{ex} %[3X15V75]
Trong trò chơi bốc thăm trúng thưởng, có 50 phiếu trong đó có 5 phiếu trúng thưởng. Bạn An rút 1 phiếu. Khả năng rút được phiếu không trúng thưởng là
\choice
{\True Có thể xảy ra}
{Chắc chắn xảy ra}
{Không thể xảy ra}
{Luôn luôn xảy ra}
\loigiai{
Có thể rút trúng phiếu trúng hoặc không trúng.\\
\textbf{Đáp án của câu hỏi là A}
}
\end{ex}

\begin{ex} %[3X15V75]
Số phiếu không trúng thưởng trong thùng là
\choice
{$40$ phiếu}
{$46$ phiếu}
{$48$ phiếu}
{\True $45$ phiếu}
\loigiai{
$50 - 5 = 45$ phiếu.\\
\textbf{Đáp án của câu hỏi là D}
}
\end{ex}

\begin{ex} %[3X15C75]
Một nhóm bạn chơi tung đồng xu 50 lần. Bạn Nam dự đoán mặt sấp xuất hiện 28 lần. Khi thực hiện, mặt sấp xuất hiện 24 lần. Số lần thực tế ít hơn dự đoán của Nam là
\choice
{$3$ lần}
{$5$ lần}
{\True $4$ lần}
{$6$ lần}
\loigiai{
$28 - 24 = 4$ lần.\\
\textbf{Đáp án của câu hỏi là C}
}
\end{ex}

\textbf{Phần 2. Trắc nghiệm đúng sai}
\setcounterpageref{ex}{0}

\begin{ex} %[3X15H75]
Xét tính đúng sai về kiểm đếm số liệu:
\choiceTF
{\True Kiểm đếm giúp thu thập số liệu một cách chính xác.}
{\True 4 nhóm (mỗi nhóm 5) tương ứng với 20 đơn vị.}
{5 nhóm (mỗi nhóm 5) và 1 gạch lẻ là 25 đơn vị.}
{\True 5 nhóm (mỗi nhóm 5) và 1 gạch lẻ là 26 đơn vị ($5 \times 5 + 1 = 26$).}
\loigiai{
\begin{itemize}
    \item Kiểm đếm chính xác: Đúng.
    \item $4 \times 5 = 20$: Đúng.
    \item $5 \times 5 + 1 = 26 \neq 25$. Ý 3 Sai, Ý 4 Đúng.
\end{itemize}
\begin{itemize}
    \item a) \textbf{Đúng}.
    \item b) \textbf{Đúng}.
    \item c) \textbf{Sai}.
    \item d) \textbf{Đúng}.
\end{itemize}
}
\end{ex}

\begin{ex} %[3X15V75]
Tung đồng xu 30 lần, ghi nhận được 18 lần mặt ngửa. Xét tính đúng sai:
\choiceTF
{\True Số lần xuất hiện mặt sấp là 12 lần ($30 - 18 = 12$).}
{\True Mặt ngửa xuất hiện nhiều hơn mặt sấp 6 lần ($18 - 12 = 6$).}
{Mặt sấp xuất hiện nhiều hơn mặt ngửa.}
{\True Mặt ngửa xuất hiện chiếm hơn một nửa tổng số lần tung.}
\loigiai{
\begin{itemize}
    \item Sấp: $30 - 18 = 12$ lần. Đúng.
    \item Hiệu: $18 - 12 = 6$ lần. Đúng.
    \item Mặt ngửa nhiều hơn (không phải mặt sấp). Sai.
    \item $18 > 15$ (hơn một nửa): Đúng.
\end{itemize}
\begin{itemize}
    \item a) \textbf{Đúng}.
    \item b) \textbf{Đúng}.
    \item c) \textbf{Sai}.
    \item d) \textbf{Đúng}.
\end{itemize}
}
\end{ex}

\textbf{Phần 3. Trả lời ngắn}
\setcounterpageref{ex}{0}

\begin{ex} %[3X15H75]
Một bạn đếm được 4 nhóm (mỗi nhóm 5 vạch) và 3 vạch rời. Tính tổng số đơn vị đếm được.
\shortans{23}
\loigiai{
\textbf{Đáp số:} $23$\\
$4 \times 5 + 3 = 23$.
}
\end{ex}

\begin{ex} %[3X15H75]
Tung đồng xu 40 lần, có 22 lần xuất hiện mặt sấp. Hỏi có bao nhiêu lần xuất hiện mặt ngửa?
\shortans{18}
\loigiai{
\textbf{Đáp số:} $18$\\
$40 - 22 = 18$ lần.
}
\end{ex}

\begin{ex} %[3X15V75]
Kiểm đếm số cây trong vườn: có 35 cây chuối, 45 cây cam và 20 cây bưởi. Tính tổng số cây trong vườn.
\shortans{100}
\loigiai{
\textbf{Đáp số:} $100$\\
$35 + 45 + 20 = 100$ cây.
}
\end{ex}

\begin{ex} %[3X15V75]
Số cây cam nhiều hơn số cây bưởi bao nhiêu cây?
\shortans{25}
\loigiai{
\textbf{Đáp số:} $25$\\
$45 - 20 = 25$ cây.
}
\end{ex}

\begin{ex} %[3X15V75]
Gieo xúc xắc 50 lần, thấy mặt 1 chấm xuất hiện 8 lần, mặt 2 chấm xuất hiện 12 lần. Cả hai mặt này xuất hiện tổng cộng bao nhiêu lần?
\shortans{20}
\loigiai{
\textbf{Đáp số:} $20$\\
$8 + 12 = 20$ lần.
}
\end{ex}

\begin{ex} %[3X15C75]
Một hộp có 30 viên bi gồm bi đỏ, xanh và vàng. Biết có 12 bi đỏ, số bi xanh có 10 viên. Hỏi có bao nhiêu viên bi vàng?
\shortans{8}
\loigiai{
\textbf{Đáp số:} $8$\\
$30 - (12 + 10) = 8$ viên bi vàng.
}
\end{ex}
